\section{Implement a Queue using a Linked List}
\label{sec:sq-queue-linked-list}

\problemstatement{Implement a queue supporting $\textit{enqueue}(x)$,
$\textit{dequeue}()$, $\textit{front}()$, and $\textit{isEmpty}()$, backed
by a singly linked list, each operation running in $O(1)$.}

\begin{patternbox}
A queue's FIFO discipline needs insertion at \textbf{one} end and removal
from the \textbf{other} -- unlike a stack, a single \textit{head} pointer
is not enough, because removing from the head and inserting at the head
would give LIFO, not FIFO. The keyword to notice: this is the first
structure in this chapter that genuinely requires \textbf{two} pointers
into the list (\textit{head} and \textit{tail}) to stay $O(1)$.
\end{patternbox}

\subsection*{Understanding the problem carefully}

If we insert new elements at the \textbf{tail} and remove from the
\textbf{head}, the earliest-inserted, not-yet-removed element is always
at the head, ready to be dequeued -- exactly FIFO order. Maintaining only
a \textit{head} pointer would make tail-insertion cost $O(n)$ (we would
have to walk the whole list to find the current last node), so we keep an
explicit \textit{tail} pointer too, updated on every enqueue.

\begin{ideabox}[Two Pointers: Head for Removal, Tail for Insertion]
$\textit{enqueue}(x)$ creates a new node and links it after the current
\textit{tail} (or sets it as both \textit{head} and \textit{tail} if the
queue was empty), then updates \textit{tail} to point to this new node.
$\textit{dequeue}()$ reads \textit{head}'s value, then advances
$\textit{head} \gets \textit{head}.\textit{next}$ (and resets
\textit{tail} to null if the queue becomes empty).
\end{ideabox}

\subsection*{Algorithm}

\begin{enumerate}
  \item $\textit{head}, \textit{tail} \gets$ null initially.
  \item $\textit{enqueue}(x)$: $\textit{node} \gets$ new node with value
        $x$. If $\textit{tail} = $ null: $\textit{head} \gets
        \textit{tail} \gets \textit{node}$. Else:
        $\textit{tail}.\textit{next} \gets \textit{node}$;
        $\textit{tail} \gets \textit{node}$.
  \item $\textit{dequeue}()$: $\textit{val} \gets
        \textit{head}.\textit{val}$; $\textit{head} \gets
        \textit{head}.\textit{next}$; if $\textit{head} = $ null:
        $\textit{tail} \gets$ null; return \textit{val}.
  \item $\textit{front}()$: return $\textit{head}.\textit{val}$.
  \item $\textit{isEmpty}()$: return $(\textit{head} = \text{null})$.
\end{enumerate}

\subsection*{Dry run}

\dryrun{Enqueue $1,2,3$; dequeue once.}

\begin{itemize}
  \item Enqueue $1$: $\textit{head}=\textit{tail}=1$.
  \item Enqueue $2$: $1 \to 2$; $\textit{head}=1$, $\textit{tail}=2$.
  \item Enqueue $3$: $1\to2\to3$; $\textit{head}=1$, $\textit{tail}=3$.
  \item Dequeue: returns $1$; $\textit{head}=2$ now; list $2\to3$.
\end{itemize}

\subsection*{C++ implementation}

\begin{lstlisting}
#include <bits/stdc++.h>
using namespace std;

class Queue {
    struct Node {
        int val;
        Node* next;
        Node(int v) : val(v), next(nullptr) {}
    };
    Node* head = nullptr;
    Node* tail = nullptr;

public:
    void enqueue(int x) {
        Node* node = new Node(x);
        if (!tail) {
            head = tail = node;
        } else {
            tail->next = node;
            tail = node;
        }
    }

    int dequeue() {
        int val = head->val;
        Node* old = head;
        head = head->next;
        if (!head) tail = nullptr;
        delete old;
        return val;
    }

    int front() {
        return head->val;
    }

    bool isEmpty() {
        return head == nullptr;
    }
};

int main() {
    Queue q;
    q.enqueue(1); q.enqueue(2); q.enqueue(3);
    cout << "Dequeued: " << q.dequeue() << endl;
    cout << "Front: " << q.front() << endl;
    return 0;
}
\end{lstlisting}

\begin{complexitybox}
\textbf{Time Complexity.} $O(1)$ for every operation, since \textit{tail}
removes the need to traverse the list on enqueue.
\textbf{Space Complexity.} $O(n)$ for $n$ stored elements.
\end{complexitybox}

\begin{takeawaybox}
The need for two pointers (\textit{head} and \textit{tail}) rather than
one is the structural signature of FIFO order on a singly linked list.
This same head/tail pairing reappears, in a different form, inside the
doubly linked list used by LRU Cache (Section~\ref{sec:sq-lru-cache}).
\end{takeawaybox}
